\subsection{The complete regular chain}

The regular primary, wedge and secondary comparisons form one chain.
A common crossing direction on their actual overlaps gives a single
zero bound for $K$. Near the parabolic corner, the primary coefficient
cone shows that the outward direction is $-\operatorname{sign}C$,
with $C=\bar d/\lambda$.

\begin{proposition}\label{rw:regular-chain}
Assume the fixed-field inequality \eqref{rw:counting:forcing}.
The regular initial part of the ordered physical interval has at
most one zero of $K$, and that zero is simple. Near $\mathcal P$
this assertion includes the finite connection, the wedge and the
regular secondary tags, with their actual overlaps.

If a primary regular zero has occurred, its subsequent sign and
outward derivative have sign $-\sigma_1$ on the remaining
regular chain. The same conclusion holds if its initial point
has
\begin{equation}\label{rw:merge:strong-entry}
 \tau_1\sigma_1<0,\qquad
 \operatorname{sign}K=-\operatorname{sign}\sigma_1;
\end{equation}
in that case the regular chain is zero-free.
If the first zero instead occurs in the wedge or in the regular
secondary part, its subsequent sign and outward derivative have
sign $-\bar d$, and $\bar a\bar d<0$.
In particular a secondary terminal after that zero receives
$\bar aC<0$ and $K$ of sign $-C$.
\end{proposition}
\begin{proof}
Saturation and coefficient transport give the positive-central
interval with fixed-fibre forcing. The common-input identity
\eqref{rw:counting:input-change} identifies $K$ on its overlaps.

On the primary compact regular part
Lemma~\ref{rw:finite:regular-direction} gives at most one
simple zero and its outward direction $-\sigma_1$.
At a zero, or under \eqref{rw:merge:strong-entry}, the
positive column quotient on a fixed regular overlap gives
\begin{equation}\label{rw:merge:primary-cone}
 |\tau_1|\le C|\sigma_1|,\qquad
 |\nu|\le C e^{-p_\sharp^2/(20\epsilon)}|\sigma_1|.
\end{equation}
The second inequality involves coefficients only, so it persists
as the input moves along the chain.
Near $\mathcal P$, $\bar B+\bar m$ has a fixed positive
margin. Thus $\bar c$ has the sign of $\sigma_1$,
$|\bar a|\le C|\bar c|$, and
$\bar a\bar c<0$: the latter follows from
$(\bar a+\bar c)\bar c<0$.
Equation \eqref{rw:merge:primary-cone} makes
$\nu=o(|\bar c|)$, so $\bar d=\bar c+\nu$ has that same
sign and $\bar a\bar d<0$.
For $\lambda>0$, $C=\bar d/\lambda$ has that sign too.
No division of the error estimate by $\lambda$ was used.

On the wedge the two columns in
\eqref{rw:secondary:wedge-columns}, after the fixed-fibre
forcing column is absorbed, have negative margins and a
strictly monotone quotient in the actual outward direction.
Hence any zero is simple, requires $\bar a\bar d<0$ and
crosses in direction $-\bar d$. The regular secondary
formula \eqref{rw:secondary:regular-limit} has the same
property: its lower constant has sign $-\bar a$, and
$\log(Q/P)$ has sign $C$, with the derivative margin of
Lemma~\ref{rw:secondary:chain-direction}.
Near the center this follows from its divided columns;
away from the center the nonzero value or derivative margin
proves it directly. A zero with $C=0$ would force the
center identity, so does not give an isolated zero.
The strict-focus reciprocal endpoint uses the stated value
or derivative estimate in that same regular secondary
preparation.

Apply Lemma~\ref{rw:counting:merge} on the actual open overlaps.
After a primary zero, the coefficient cone makes $\bar d$ and
$\sigma_1$ have the same sign, so all later crossings have the
same direction. Without a primary zero, all wedge and secondary
crossings have direction $-\bar d$. This gives one zero on the
union; the local direction lemmas keep $K$ departing once it has
sign $-\bar d$.

The finite lower-double connections in
Lemma~\ref{rw:fibre:parabolic-connections} require a separate
observation. On the lower-cone boundary their regular prefix,
finite bridge, wedge and every retained compact regular
secondary interval have one strict action sign, opposite to
$\bar a$. Their lower action dominates on every retained compact
secondary part, so they contribute no zero; a singular end remains
terminal. On the primary prefix
$\tau_1$ has the sign of $\bar a$. If
$\tau_1\sigma_1<0$, the sign $-\sigma_1$ is therefore
opposite to this actual value of $K$. Thus the strengthened
entrance cannot occur in these boundary neighborhoods.

Under \eqref{rw:merge:strong-entry}, the primary part departs from
zero. Its cone \eqref{rw:merge:primary-cone} carries the same direction
to the later pieces, excluding a first zero there. This proves the
strengthened assertion.
\end{proof}
